\chapter{The real region and the bands}\label{chap:region}

When the choice of the sieve (Chapter~\ref{chap:sieve-numbers}) takes the tube or the real region,
the main search sieves the rows $(a, b)$ of a first coordinate~$a$ by \texttt{sift\_bands}
(\csource{j-sift.c}{815}{878}) instead of row by row: the rows are taken in bands of consecutive
rows, and in each row only some ranges of words are sieved. In the tube these are the ranges that
the analysis of the tube collects (Chapter~\ref{chap:tube}), possibly cut to the real region; in
the real region they are the ranges that \texttt{row\_region} (\csource{j-sift.c}{484}{556})
leaves of the row once it has taken out the third coordinates that no point can have. This chapter
says what \texttt{row\_region} keeps (Section~\ref{sec:rowRegion}), why the ends of the intervals
that it takes out, computed in doubles, lie on the right side of the exact ones when the height is
at most~$2^{26}$ (Section~\ref{sec:rowRegionEnds}), why no point lies in them
(Section~\ref{sec:realRegion}), what \texttt{sift\_bands} prints (Section~\ref{sec:bands}), and
that the test of the bands by which the filtered enumeration of Chapter~\ref{chap:theorems}
describes it keeps every point (Section~\ref{sec:bandsSound}).

Throughout, the third coordinates~$c$ of a row lie in the words $\lfloor c/64 \rfloor$, and the
words of the row are those from $w_{\mathrm{low}}$ to $w_{\mathrm{high}}$, excluded (the statics
\texttt{w\_low\_all} and \texttt{w\_high\_all}). A double~$x$ has the value
$v(x) \in \QQ \cup \{\pm\infty\}$ (Definition~\ref{def:doubleValue}); where no confusion can arise
we write~$x$ for the value of a finite double as well.

\section{The real region of a row}\label{sec:rowRegion}

The intervals~$\mathrm{neg}$ that \texttt{real\_roots\_init} finds (Chapter~\ref{chap:doubles}) are
at most four, each given by its shrunk ends $n_l < n_h$, between which the polynomial~$f$ of the
curve searched is negative (Definition~\ref{def:NegativeIntervals.Sound}). They exclude third
coordinates: if a point lies over $(a, b, c)$ and $a \ne 0$, then~$f$ takes a value of at
least~$0$ at every real root of $a t^2 - b t + c$ (Section~\ref{sec:realRegion}), so no real root
lies inside such an interval. At a root~$t$ the third coordinate is $c = q(t) = bt - at^2$. For
$a > 0$, $q$ increases up to the vertex $m = b/(2a)$ and decreases after it; so the third
coordinates with the larger root in $(n_l, n_h)$ form the interval from $q(n_h)$ to $q(l)$ for
$l = \max(n_l, m)$, and those with the smaller root in it the interval from $q(n_l)$ to $q(h)$ for
$h = \min(n_h, m)$. For $a = 0$ the only root is $c/b$, and $b > 0$ in the rows of the search: the
third coordinates excluded are those between $b n_l$ and $b n_h$. The function
\texttt{row\_region} computes these intervals in doubles, makes each smaller by a margin for the
rounding, and takes out of the row the words that lie wholly inside one of them.

\begin{definition}[\texttt{row\_region}]\label{def:rowRegion}
  \lean{JPoints.rowRegion}
  \leanok
  \uses{def:Double.add, def:Double.ofDyadic}
  For $a \ge 0$ (and $b > 0$ when $a = 0$), $\mathrm{rowRegion}(\mathrm{neg}, \mathrm{sd}, a, b)$ is
  the real region of the row $(a, b)$ as \texttt{row\_region} computes it: ranges of words
  $[\mathrm{lo}_k, \mathrm{hi}_k)$ of the row, in order, and their number. Without an interval on
  which~$f$ is negative it is the whole row $[w_{\mathrm{low}}, w_{\mathrm{high}})$. Otherwise let
  $A$ and~$B$ be the doubles of~$a$ and~$b$, let $\widehat{q}(x)$ be the double $B x - A x x$,
  computed in this order, and let $(n_l, n_h)$ run through the shrunk intervals of~$\mathrm{neg}$.
  \begin{enumerate}
    \item The excluded intervals of the third coordinate: for $a > 0$ and the vertex
      $m = 0.5 B / A$, on the branch of the larger root, when $l = \max(n_l, m) < n_h$, the interval
      from $\widehat{q}(n_h)$ ($-\infty$ for $n_h = \infty$) to $\widehat{q}(l)$, and on the branch
      of the smaller root, when $n_l < h = \min(n_h, m)$, the interval from $\widehat{q}(n_l)$
      ($-\infty$ for $n_l = -\infty$) to $\widehat{q}(h)$; for $a = 0$, the interval from $B n_l$
      to $B n_h$, an infinite end kept.
    \item Each lower end $x > -\infty$ raised to $x + \mu(x)$ and each upper end $x < \infty$
      lowered to $x - \mu(x)$, for the margin $\mu(x) = 10^{-9}(|x| + 1) + 1$, in doubles.
    \item The intervals sorted by their lower ends, by insertion.
    \item From $w = w_{\mathrm{low}}$ on, for each interval $(x_l, x_h)$ in turn, unless
      $x_l \ge \mathtt{ROW\_BEYOND}$, $x_h \le -\mathtt{ROW\_BEYOND}$ or $x_h - x_l < 64$: the words
      $[w_l, w_h)$ wholly inside it, where $w_l = \lceil ((\mathrm{long})\,x_l + 2)/64 \rceil$, or
      $w_{\mathrm{low}}$ when $x_l \le -\mathtt{ROW\_BEYOND}$, and
      $w_h = \lfloor ((\mathrm{long})\,x_h - 1)/64 \rfloor$, or $w_{\mathrm{high}}$ when
      $x_h \ge \mathtt{ROW\_BEYOND}$, cut to $[w, w_{\mathrm{high}})$; when these are not empty,
      the range $[w, w_l)$ kept if it is not empty, and~$w$ moved to~$w_h$. At the end the range
      $[w, w_{\mathrm{high}})$ is kept if it is not empty.
  \end{enumerate}
\end{definition}

Here $\mathtt{ROW\_BEYOND} = 4 \cdot 10^9$ lies beyond every third coordinate of a row, and
$(\mathrm{long})$ truncates towards~$0$; so a word from~$w_l$ on holds only third coordinates
$c > x_l + 1$, and a word below~$w_h$ only $c < x_h - 1$ (\lname{JPoints.le\_or\_lt\_of\_rowRegionEnds\_fst\_le},
\lname{JPoints.le\_or\_lt\_of\_lt\_rowRegionEnds\_snd}).

\formalnote{The model computes with the bits of the doubles, each operation as gcc and glibc
compute it, in the order of the C code, with every comparison false when either side is a NaN. It
returns the ranges as arrays of their number, for the arrays of \texttt{ROW\_RANGES} entries of the
caller, of which at most $2 \cdot 4 + 1$ are used, and the excluded intervals as arrays, for the
local arrays of eight doubles.}

What \texttt{row\_region} takes out is described by the words of item~(4) alone, before the cut to
$[w, w_{\mathrm{high}})$ and without the test of the width: a description of more words than it
takes out, which is the direction that the soundness of the real region needs.

\begin{definition}[What \texttt{row\_region} excludes]\label{def:RegionExcludes}
  \lean{JPoints.rowRegionValue, JPoints.RegionExcludes, JPoints.RegionCovers}
  \leanok
  \uses{def:Double.add, def:Double.ofDyadic}
  For the row $(a, b)$ and a double~$x$, the double $\widehat{q}(x) = (\mathrm{double})\,b \cdot x -
  (\mathrm{double})\,a \cdot x \cdot x$, computed in this order, is
  $\mathrm{rowRegionValue}(a, b, x)$: an end of an excluded interval. The word~$v$ is \emph{taken out} for the excluded interval
  $(x_l, x_h)$, $\mathrm{RegionExcludes}(\mathrm{sd}, x_l, x_h, v)$, when the interval does not lie
  beyond the row (the comparisons $\mathtt{ROW\_BEYOND} \le x_l$ and
  $x_h \le -\mathtt{ROW\_BEYOND}$ are false) and $w_l \le v < w_h$ for the words $[w_l, w_h)$ of
  item~(4) of Definition~\ref{def:rowRegion}. Ranges~$r$ \emph{cover} the row below~$w$ for the
  first~$k$ intervals of the arrays $x_{\mathrm{lo}}$ and~$x_{\mathrm{hi}}$ ($\mathrm{RegionCovers}$)
  when every word
  $w_{\mathrm{low}} \le v < w$ lies in a range of~$r$ or is taken out for one of the intervals
  $(x_{\mathrm{lo}}[i], x_{\mathrm{hi}}[i])$, $i < k$.
\end{definition}

\begin{lemma}[\texttt{row\_region} covers the row]\label{lem:rowRegion_cover}
  \lean{JPoints.rowRegion_cover}
  \leanok
  \uses{def:rowRegion, def:RegionExcludes, def:Double.add, def:Double.ofDyadic}
  Suppose that $w_{\mathrm{low}} < w_{\mathrm{high}}$. Then every word $w_{\mathrm{low}} \le v <
  w_{\mathrm{high}}$ of the row lies in a range of $\mathrm{rowRegion}(\mathrm{neg}, \mathrm{sd}, a,
  b)$, or there are an index $i < \mathrm{numNeg}$ and an excluded interval $(\mathrm{lo},
  \mathrm{hi})$ that \texttt{row\_region} takes from the~$i$-th interval of~$\mathrm{neg}$ for the
  row $(a, b)$ (item~(1) of Definition~\ref{def:rowRegion}, \lname{JPoints.RegionInterval}) such
  that~$v$ is taken out for that interval raised and lowered by its margins as in item~(2).
\end{lemma}

\begin{proof}
  \leanok
  By the contracts of the parts of \texttt{row\_region}, each proved by \texttt{vcgen} with a named
  invariant: the excluded intervals come from the intervals of~$\mathrm{neg}$, by the two branches
  for $a > 0$ and by $B n_l$ and $B n_h$ otherwise (\lname{JPoints.rowRegionPositive\_spec},
  \lname{JPoints.rowRegionZero\_spec}); the margins move each end
  (\lname{JPoints.rowRegionMargins\_spec}); the insertion sort only permutes the pairs of ends
  (\lname{JPoints.rowRegionSort\_spec}); and the placing of the ranges keeps the ranges and the
  words taken out covering the row below~$w$ (\lname{JPoints.rowRegionWords\_cover}): the words
  from~$w$ to~$w_l$ go into the range $[w, w_l)$, those from~$w_l$ to~$w_h$ are taken out, and the
  cut to $[w, w_{\mathrm{high}})$ removes only words that are covered already or lie beyond the
  row. Without intervals the single range is the whole row.
\end{proof}

\begin{lemma}[The ranges of \texttt{row\_region}]\label{lem:rowRegion_spec}
  \lean{JPoints.rowRegion_spec}
  \leanok
  \uses{def:rowRegion}
  Suppose that $w_{\mathrm{low}} < w_{\mathrm{high}}$. Then the ranges of
  $\mathrm{rowRegion}(\mathrm{neg}, \mathrm{sd}, a, b)$ lie in $[w_{\mathrm{low}},
  w_{\mathrm{high}})$ in increasing order ($\mathrm{RangesIn}$): each is nonempty, and each ends
  before the next begins.
\end{lemma}

\begin{proof}
  \leanok
  \uses{def:Double.add, def:Double.ofDyadic}
  An invariant of the placing of the ranges: $w_{\mathrm{low}} \le w \le w_{\mathrm{high}}$, and the
  ranges placed so far lie below~$w$ in increasing order; a range $[w, w_l)$ is placed only when it
  is nonempty, and~$w$ moves to $w_h > w_l$.
\end{proof}

\section{The ends in doubles and the height bound}\label{sec:rowRegionEnds}

For the argument of Section~\ref{sec:realRegion}, a word that \texttt{row\_region} takes out must
hold only third coordinates~$c$ that lie strictly between the \emph{exact} values of~$q$ at the
ends, such as $q(n_h) < c < q(l)$ on the branch of the larger root. By the slack of the words, such
a~$c$ satisfies $x_l + 1 < c < x_h - 1$ for the raised and lowered ends; the margin~$\mu$ and this
slack must cover the error of~$\widehat{q}$. Where $|q(x)|$ is comparable to the products $bx$
and $ax^2$, the error of $\widehat{q}(x)$ is a small multiple of $|q(x)|/2^{53}$, far below the
relative margin $10^{-9}$. The delicate place is near $x = b/a$: there $q(x) = x(b - ax)$ vanishes,
the products, both close to $b^2/a$, cancel, the error of $\widehat{q}(x)$ is a matter of the last
places of the products, and the margin is about~$1$. For $|b| \le 2^{26}$ the products there stay
below~$2^{53}$, where every rounding errs by at most half a unit in the last place, at
most~$1/2$, and the error stays below $7/4$, which the margin and the slack cover together.

The bound $|b| \le 2^{26}$ comes from the height bound: \texttt{read\_input} refuses height bounds
above \texttt{MAX\_SEARCH\_HEIGHT}, which is~$2^{26}$ (\lname{JPoints.maxSearchHeight}), because at
larger heights the products near $x = b/a$ could exceed~$2^{53}$, where the doubles lie~$2$ and more
apart, and their rounding errors, which grow with them, would not be covered by a margin of
about~$1$: an excluded interval could then reach beyond the exact one and take out the word of a
point.

\begin{lemma}[The error of $B x - A x x$]\label{lem:abs_sub_le_of_roundings}
  \lean{JPoints.abs_sub_le_of_roundings}
  \leanok
  Let $a \ge 1$, $b$ and~$x$ be rationals with $b^2 \le 2^{52}$ and $|x| \le 2^{480}$, and let
  $r$, $p_1$, $p_2$ and~$e$ be rationals within the error bounds of the roundings to nearest of $ax$,
  $bx$, $rx$ and $p_1 - p_2$:
  \[ |r - ax| \le \frac{|ax|}{2^{53}} + 2^{-1075}\,, \qquad
     |p_1 - bx| \le \frac{|bx|}{2^{53}} + 2^{-1075}\,, \]
  \[ |p_2 - rx| \le \frac{|rx|}{2^{53}} + 2^{-1075}\,, \qquad
     |e - (p_1 - p_2)| \le \frac{|p_1 - p_2|}{2^{53}} + 2^{-1075}\,. \]
  Then $q = bx - ax^2$ satisfies $|e - q| \le 7/4 + |q|/2^{40}$.
\end{lemma}

\begin{proof}
  \leanok
  The error of~$e$ is at most $|q|/2^{53}$ for the last rounding, plus $1 + 2^{-53}$ times the
  error of $p_1 - p_2$, which the products $P = |b|\,|x|$ and $Q = ax^2$ bound. Near $x = b/a$,
  when $16 a |x| \le 17 |b|$, the products satisfy $P \le 17 b^2/16$ and $Q \le 289 b^2/256$, so
  they lie below~$2^{53}$, and the three errors of the products are at most about~$0.57$ each.
  Away from it, $P \le 16 Q/17$ and $Q \le 17 |q|$, and the error is relative. The terms
  $2^{-1075}$, of the subnormal range, are negligible for $|x| \le 2^{480}$.
\end{proof}

\begin{lemma}[The margins of the ends]\label{lem:Double.value_lt_of_marginLo}
  \lean{JPoints.Double.value_lt_of_marginLo, JPoints.Double.lt_value_of_marginHi}
  \leanok
  \uses{def:Double.toRat, def:Double.add, def:Double.ofDyadic}
  Let $1 \le a \le 2^{26}$ and $|b| \le 2^{26}$, let~$x$ be a finite double, $e = \widehat{q}(x)$,
  and $q = bx - ax^2$ for the value~$x$ of the double.
  \begin{enumerate}
    \item Let $x_l = e + \mu(e)$ when $-\infty < e$ and $x_l = e$ otherwise, and let
      $c \ge -2^{26}$ be an integer. If~$x_l$ is a NaN, or $x_l \le -\mathtt{ROW\_BEYOND}$, or
      $x_l$ is finite with $x_l + 1 < c$, then $q < c$.
    \item Let $x_h = e - \mu(e)$ when $e < \infty$ and $x_h = e$ otherwise, and let
      $c \le 2^{26}$ be an integer. If~$x_h$ is no NaN, $x_h \le -\mathtt{ROW\_BEYOND}$ is false,
      and $x_h \ge \mathtt{ROW\_BEYOND}$ or~$x_h$ is finite with $c < x_h - 1$, then $c < q$.
  \end{enumerate}
\end{lemma}

\begin{proof}
  \leanok
  \uses{def:doubleValue, lem:Double.lt_eq_true_iff, def:Double.ofDyadic, lem:Double.ofDyadic_exact, def:Double.add, def:Double.IsRounding, lem:Double.isRounding_add, lem:abs_sub_le_of_roundings, lem:Double.isNaN_or_doubleValue_value_le}
  The doubles of~$a$ and~$b$ are exact. For $|x| \le 2^{480}$ the operations of $\widehat{q}(x)$
  stay finite and round to nearest, and $|e - q| \le 7/4 + |q|/2^{40}$
  by Lemma~\ref{lem:abs_sub_le_of_roundings};
  the margin $\mu(e)$, computed with three roundings of its own, is finite and at least
  $1 + 10^{-9}|e|$ up to a negligible error, and together with the slack~$1$ it exceeds the error.
  An end $x_l \le -\mathtt{ROW\_BEYOND}$ puts~$q$ below $-\mathtt{ROW\_BEYOND} + 2 < -2^{26}$ in
  the same way. For $|x| > 2^{480}$, $q \le |bx| - ax^2 < -2^{26} \le c$ whatever the double~$e$;
  and there~$e$ is a NaN or at most $-2^{900}$
  (Lemma~\ref{lem:Double.isNaN_or_doubleValue_value_le}), so that~$x_h$ is a NaN or at most
  $-\mathtt{ROW\_BEYOND}$, and the hypothesis of item~(2) fails.
\end{proof}

\begin{lemma}[The ends far out]\label{lem:Double.isNaN_or_doubleValue_value_le}
  \lean{JPoints.Double.isNaN_or_doubleValue_value_le}
  \leanok
  \uses{def:Double.toRat, def:doubleValue, def:Double.add, def:Double.ofDyadic}
  For $1 \le a \le 2^{26}$, $|b| \le 2^{26}$ and a finite double~$x$ with $|x| > 2^{480}$, the
  double $\widehat{q}(x)$ is a NaN or has a value of at most $-2^{900}$.
\end{lemma}

\begin{proof}
  \leanok
  \uses{def:Double.ofDyadic, lem:Double.ofDyadic_exact, def:Double.add, def:Double.IsRounding, lem:Double.isRounding_add}
  The product $(A x) x$ is $+\infty$ or finite and at least $x^2/4$, while $B x$ is an infinity or
  at most about $2^{26}|x|$ in absolute value; the rounded difference is then $-\infty$, a NaN, or
  at most $-2^{900}$.
\end{proof}

\begin{lemma}[The ends for $a = 0$]\label{lem:Double.lt_div_of_marginLo}
  \lean{JPoints.Double.lt_div_of_marginLo, JPoints.Double.div_lt_of_marginHi}
  \leanok
  \uses{def:Double.toRat, def:doubleValue, def:Double.add, def:Double.ofDyadic}
  Let $1 \le b \le 2^{26}$ and let~$n$ be a double that is no NaN.
  \begin{enumerate}
    \item Let $\mathrm{lo} = -\infty$ when $n = -\infty$ and $\mathrm{lo} = B n$ otherwise, let
      $x_l$ be~$\mathrm{lo}$ raised by its margin as in item~(2) of Definition~\ref{def:rowRegion},
      and let $c \ge -2^{26}$ be an integer. If $x_l \ge \mathtt{ROW\_BEYOND}$ is false and, unless
      $x_l$ is a NaN, $x_l \le -\mathtt{ROW\_BEYOND}$ or~$x_l$ is finite with $x_l + 1 < c$, then
      $v(n) < c/b$.
    \item Let $\mathrm{hi} = \infty$ when $n = \infty$ and $\mathrm{hi} = B n$ otherwise, let~$x_h$
      be~$\mathrm{hi}$ lowered by its margin, and let $c \le 2^{26}$ be an integer. If~$x_h$ is no
      NaN, $x_h \le -\mathtt{ROW\_BEYOND}$ is false, and $x_h \ge \mathtt{ROW\_BEYOND}$ or~$x_h$ is
      finite with $c < x_h - 1$, then $c/b < v(n)$.
  \end{enumerate}
\end{lemma}

\begin{proof}
  \leanok
  \uses{lem:Double.lt_eq_true_iff, def:Double.ofDyadic, lem:Double.ofDyadic_exact, def:Double.add, def:Double.IsRounding, lem:Double.isRounding_add}
  An end $n = -\infty$ in item~(1), or $n = \infty$ in item~(2), gives the claim at once; an
  infinite~$n$ of the other sign makes~$B n$ infinite and $x_l \ge \mathtt{ROW\_BEYOND}$ (or
  $x_h \le -\mathtt{ROW\_BEYOND}$), which is excluded. For a finite~$n$ with $|n| \le 2^{900}$ the product $B n$ is finite and within
  $7/4 + |bn|/2^{40}$ of~$bn$, and the margin and the slack cover the error as in
  Lemma~\ref{lem:Double.value_lt_of_marginLo}. Far out, $n > 2^{900}$ makes~$x_l$ at least
  $\mathtt{ROW\_BEYOND}$ and $n < -2^{900}$ puts~$n$ below $c/b$ (dually for item~(2)).
\end{proof}

\section{The curve over the real region}\label{sec:realRegion}

A point over a triple $(a, b, c)$ that comes from a rational point of the Jacobian gives real
points of the curve over the real roots of $a t^2 - b t + c$: the lifting test of the
specification provides rationals $\alpha$ and~$\beta$ with $a^4 f(x) = (\alpha x + \beta)^2$ at
every root~$x$, and for $a = 0$ it says that $b^6 f(c/b)$ is a square. Inside an interval
of~$\mathrm{neg}$, on the other hand, $f$ is negative. An excluded interval that took out the word
of such a~$c$ would therefore, by the bounds of Section~\ref{sec:rowRegionEnds} and the
intermediate value theorem, put a root inside an interval of~$\mathrm{neg}$: a contradiction.

\begin{lemma}[$f$ at the real roots over a point]\label{lem:sum_nonneg_of_hasPointOver}
  \lean{JPoints.sum_nonneg_of_hasPointOver, JPoints.eval_poly_div_nonneg_of_hasPointOver}
  \leanok
  \uses{def:poly, def:HasPointOver}
  Let $f = \mathrm{poly}(\mathrm{cs})$ be the polynomial of the coefficient list
  $\mathrm{cs} = (f_0, f_1, \dots)$.
  \begin{enumerate}
    \item If $a \ne 0$, a point lies over $(a, b, c)$, and $x \in \mathbb{R}$ is a root of
      $a t^2 - b t + c$, then $\sum_{k < 7} f_k x^k \ge 0$.
    \item If~$\mathrm{cs}$ has at most seven entries, $b \ne 0$, and a point lies over
      $(0, b, c)$, then $f(c/b) \ge 0$.
  \end{enumerate}
\end{lemma}

\begin{proof}
  \leanok
  \uses{def:kummer, def:auxA2, def:Lifts, def:kummerRoots, lem:mem_kummerRoots_iff}
  A point over the triple is a point $(ka, kb, kc, d)$ of the Kummer surface, $k \ne 0$, that lifts
  to the Jacobian. For $a \ne 0$ the lifting test gives rationals $\alpha$ and~$\beta$ with
  $\alpha^2 = A_2$, $\beta^2 = B_2$ and $2\alpha\beta$ the third polynomial of the test, and at a
  root~$x$ the remainder of $(ka)^4 f - A_2 t^2 - 2\alpha\beta\, t - B_2$ modulo the quadratic
  vanishes, so $(ka)^4 f(x) = (\alpha x + \beta)^2 \ge 0$. For $a = 0$ the test says that
  $F(kc, kb) = k^6 b^6 f(c/b)$ is a square, for the sextic form~$F$ of~$f$.
\end{proof}

\begin{definition}[Inside a negative interval]\label{def:NegativeIntervals.Contains}
  \lean{JPoints.NegativeIntervals.Contains}
  \leanok
  \uses{def:doubleValue}
  The real number~$x$ lies inside the~$i$-th interval of~$\mathrm{neg}$,
  $\mathrm{Contains}(\mathrm{neg}, i, x)$, when there are reals $L < x < U$ such that every
  rational~$q$ with $L < q < U$ lies strictly between the values of the shrunk ends of the
  interval.
\end{definition}

\begin{lemma}[$f$ is negative inside a sound interval]\label{lem:NegativeIntervals.Sound.eval_map_neg_of_contains}
  \lean{JPoints.NegativeIntervals.Sound.eval_map_neg_of_contains}
  \leanok
  \uses{def:poly, def:NegativeIntervals.Sound, def:NegativeIntervals.Contains}
  If the intervals~$\mathrm{neg}$ are sound for~$\mathrm{cs}$, $f = \mathrm{poly}(\mathrm{cs})$ is
  squarefree, $i < \mathrm{numNeg}$ and~$x$ lies inside the~$i$-th interval, then $f(x) < 0$ for~$f$
  over~$\mathbb{R}$.
\end{lemma}

\begin{proof}
  \leanok
  \uses{def:doubleValue}
  Soundness makes~$f$ negative at the rationals between the values of the shrunk ends, hence at the
  rationals of $(L, U)$. A squarefree polynomial over~$\QQ$ that is negative at the rational points
  of an open interval of~$\mathbb{R}$ is at most~$0$ on it by continuity, and a real root inside
  would be a local maximum, where the derivative vanishes, which squarefreeness rules out.
\end{proof}

\begin{lemma}[An excluded word lies inside a negative interval]\label{lem:NegativeIntervals.Sound.exists_contains_of_regionExcludes_from}
  \lean{JPoints.NegativeIntervals.Sound.exists_contains_of_regionExcludes_from, JPoints.NegativeIntervals.Sound.exists_contains_of_regionExcludes_upTo}
  \leanok
  \uses{def:NegativeIntervals.Sound, def:RegionExcludes, def:NegativeIntervals.Contains, def:Double.add, def:Double.ofDyadic}
  Let the intervals~$\mathrm{neg}$ be sound for~$\mathrm{cs}$, let $i < \mathrm{numNeg}$, with the
  shrunk ends~$n_l$ and~$n_h$, let $1 \le a \le 2^{26}$ and $|b|, |c| \le 2^{26}$, and let~$m$, $l$
  and~$h$ be as in item~(1) of Definition~\ref{def:rowRegion}.
  \begin{enumerate}
    \item If $l < n_h$ and the word $\lfloor c/64 \rfloor$ is taken out for the excluded interval
      from $\widehat{q}(n_h)$ ($-\infty$ when $n_h = \infty$) to $\widehat{q}(l)$, raised and
      lowered by its margins, then $a t^2 - b t + c$ has a real root~$t$ inside the~$i$-th interval.
    \item The same holds when $n_l < h$ and the word is taken out for the excluded interval from
      $\widehat{q}(n_l)$ ($-\infty$ when $n_l = -\infty$) to $\widehat{q}(h)$, raised and lowered by
      its margins.
  \end{enumerate}
\end{lemma}

\begin{proof}
  \leanok
  \uses{def:Double.toRat, def:doubleValue, lem:Double.lt_eq_true_iff, lem:Double.ofDyadic_exact, def:Double.IsRounding, lem:Double.isRounding_add, lem:Double.value_lt_of_marginLo}
  The word lies wholly inside the raised and lowered interval with the slack~$1$, so
  Lemma~\ref{lem:Double.value_lt_of_marginLo} gives the exact inequalities $q(n_h) < c < q(l)$ in
  item~(1) (only $c < q(l)$ when $n_h = \infty$) and $q(n_l) < c < q(h)$ in item~(2). As
  $a t^2 - b t + c = c - q(t)$, the quadratic is negative at~$l$ and positive at~$n_h$, or for
  large~$t$, and by the intermediate value theorem it has a root between them; likewise between
  $n_l$, or small~$t$, and~$h$. Since $n_l \le l$ and $h \le n_h$, every rational between the two
  points lies strictly between the values of the ends of the interval.
\end{proof}

\begin{lemma}[\texttt{row\_region} excludes no point]\label{lem:NegativeIntervals.Sound.not_regionExcludes_of_hasPointOver}
  \lean{JPoints.NegativeIntervals.Sound.not_regionExcludes_of_hasPointOver}
  \leanok
  \uses{def:poly, def:HasPointOver, def:NegativeIntervals.Sound, def:RegionExcludes, def:Double.add, def:Double.ofDyadic}
  Let the intervals~$\mathrm{neg}$ be sound for a list~$\mathrm{cs}$ of at most seven coefficients
  with $\mathrm{poly}(\mathrm{cs})$ squarefree, let $0 \le a \le 2^{26}$, $|b|, |c| \le 2^{26}$,
  $b > 0$ when $a = 0$, and let a point lie over $(a, b, c)$. Then for $i < \mathrm{numNeg}$ and
  every excluded interval that \texttt{row\_region} takes from the~$i$-th interval of~$\mathrm{neg}$
  for the row $(a, b)$, the word $\lfloor c/64 \rfloor$ is not taken out for that interval raised
  and lowered by its margins.
\end{lemma}

\begin{proof}
  \leanok
  \uses{def:Double.toRat, def:doubleValue, lem:Double.lt_eq_true_iff, lem:Double.lt_div_of_marginLo, lem:sum_nonneg_of_hasPointOver, def:NegativeIntervals.Contains, lem:NegativeIntervals.Sound.eval_map_neg_of_contains, lem:NegativeIntervals.Sound.exists_contains_of_regionExcludes_from}
  For $a > 0$, Lemma~\ref{lem:NegativeIntervals.Sound.exists_contains_of_regionExcludes_from}
  gives a real root inside the interval, where $f < 0$ by
  Lemma~\ref{lem:NegativeIntervals.Sound.eval_map_neg_of_contains} and $f \ge 0$ by
  Lemma~\ref{lem:sum_nonneg_of_hasPointOver}. For $a = 0$, Lemma~\ref{lem:Double.lt_div_of_marginLo}
  puts $c/b$ strictly between the values of the ends, where $f(c/b) < 0$ by soundness, against
  $f(c/b) \ge 0$.
\end{proof}

\section{The sieve by bands}\label{sec:bands}

The rows of a first coordinate~$a$ with~$b$ from $b_{\mathrm{first}}$ on by the step~$m$ are cut
into bands: a band is a block of at most \texttt{TUBE\_BAND} consecutive rows that starts on a row
of the step, the last one cut at the height bound, and its rows of the step are sieved together.

\begin{definition}[The bands of the rows]\label{def:bandOf}
  \lean{JPoints.bandOf, JPoints.bandLen, JPoints.bandStep, JPoints.bandEnd}
  \leanok
  For $m \ge 1$, $\mathrm{bandLen}(m)$ is~$1$ when $m > \mathtt{TUBE\_BAND}$ and otherwise
  $\mathtt{TUBE\_BAND} - (\mathtt{TUBE\_BAND} \bmod m)$, the largest multiple of~$m$ up to
  \texttt{TUBE\_BAND}; $\mathrm{bandStep}(m)$ is~$m$ when $m > \mathtt{TUBE\_BAND}$ and
  $\mathrm{bandLen}(m)$ otherwise; $\mathrm{bandEnd}(H, b_{\mathrm{lo}}, m) = \min(b_{\mathrm{lo}} +
  \mathrm{bandLen}(m) - 1, H)$; and $\mathrm{bandOf}(H, b_{\mathrm{first}}, m, b) = (b_{\mathrm{lo}},
  b_{\mathrm{hi}})$ is the band of \texttt{sift\_bands} that holds the row~$b$: its first row
  $b_{\mathrm{lo}} = b_{\mathrm{first}} + \lfloor (b - b_{\mathrm{first}})/\sigma \rfloor\, \sigma$
  for the step $\sigma = \mathrm{bandStep}(m)$ of the bands, and its last row $b_{\mathrm{hi}} =
  \mathrm{bandEnd}(H, b_{\mathrm{lo}}, m)$.
\end{definition}

\begin{definition}[\texttt{sift\_bands}]\label{def:siftBands}
  \lean{JPoints.tubeWord, JPoints.siftBandsTube, JPoints.siftBandsRegion, JPoints.siftBandsRow, JPoints.siftBands}
  \leanok
  \uses{def:sift, def:tubeCell, def:CheckGlobals, def:checkOnePoint, def:rowRegion}
  $\mathrm{siftBands}(\mathrm{bg}, \mathrm{sd}, s, a, b_{\mathrm{first}}, m)$ is
  \texttt{sift\_bands}, the sieve of the rows $(a, b)$ with~$b$ from $b_{\mathrm{first}}$ on by the
  step $m \ge 1$ up to the height bound~$H$, in the tube (\texttt{sieve\_mode} equal to~$1$) or in
  the real region, in the state~$s$ of the thread: the polynomials of the analysis of the tube for~$a$
  (\texttt{tube\_a\_init}); then the bands of Definition~\ref{def:bandOf} in turn, each set up, in
  the tube with the words that the analysis of the band collects (\texttt{tube\_cell} on its
  cells), and its rows of the step sieved in order ($\mathrm{siftBandsRow}$). A row is skipped in
  the tube when the analysis left it no word, and, when the condition at~$2$ is used, when $a \ne 0$
  and the class $(a \bmod 64, b \bmod 64)$ admits no third coordinate; otherwise its ranges of words
  (Definition~\ref{def:bandRowRanges}) are sieved in turn after \texttt{row\_setup}: in the tube word
  by word ($\mathrm{siftBandsTube}$), each word and-ed with the words of the sieve tables of all the
  primes of the first two stages at once and its bits handed to the third stage and the exact check
  (\texttt{tube\_word}, $\mathrm{tubeWord}$); in the region chunk by chunk, by the passes of the box
  sieve over the part of the range in each chunk ($\mathrm{siftBandsRegion}$, \texttt{sift\_range}).
  Each stops at the first point found under~\texttt{-1}; the result says whether it stopped so, with
  the state.
\end{definition}

\formalnote{The state of the thread is taken and returned, for the thread-local globals, and the
globals and the statics that the functions read are passed. The loops over the bands and over the
rows of a band, which step by $\mathrm{bandStep}(m)$ and by~$m$, are loops over the counts of their
steps, and C's \texttt{\%} is the truncating remainder.}

A row of a band in the tube with the cut to the real region (\texttt{tube\_cut}) keeps the words
that lie both in a range of the tube and in a range of the real region; \texttt{row\_cut}
(\csource{j-sift.c}{714}{724}) intersects the two lists of ranges as follows.

\begin{definition}[The cut of the tube to the real region]\label{def:cutRanges}
  \lean{JPoints.cutRanges}
  \leanok
  For lists of ranges of words with the first ranges $r = [r_1, r_2)$ and $g = [g_1, g_2)$, the cut
  is the intersection $[\max(r_1, g_1), \min(r_2, g_2))$ of the first two ranges, when it is not
  empty, followed by
  \[ \begin{cases}
       \mathrm{cutRanges}(\mathit{rs}, g :: \mathit{gs}) & \text{if } r_2 < g_2\,, \\
       \mathrm{cutRanges}(r :: \mathit{rs}, \mathit{gs}) & \text{otherwise,}
     \end{cases} \]
  for the remaining lists $\mathit{rs}$ and~$\mathit{gs}$: the range that ends first is left
  behind, that of the second list when both end together. When either list is empty, so is the cut.
\end{definition}

\begin{definition}[The ranges of a row of a band]\label{def:bandRowRanges}
  \lean{JPoints.bandRowRanges, JPoints.rangesPoints}
  \leanok
  \uses{def:rowRangePoints, def:rowRegion, def:cutRanges}
  For the ranges~$\mathit{rs}$ of the row $(a, b)$ in the bands of the analysis of the tube,
  $\mathrm{bandRowRanges}(g, \mathrm{sd}, \mathit{rs}, a, b, \mathrm{tube})$ are the ranges that
  \texttt{sift\_bands} sieves in the row: $\mathit{rs}$ in the tube without the cut;
  $\mathrm{cutRanges}(\mathit{rs}, \mathit{gs})$ for the ranges~$\mathit{gs}$ of
  $\mathrm{rowRegion}(g.\mathrm{neg}, \mathrm{sd}, a, b)$ in the tube with the cut; and the ranges
  of $\mathrm{rowRegion}(g.\mathrm{neg}, \mathrm{sd}, a, b)$ in the real region. For ranges~$R$ of
  the row, $\mathrm{rangesPoints}(g, \mathrm{sd}, a, b, R)$ lists, range after range, the points
  that $\mathrm{rowRangePoints}$ gives for the ranges $[\mathrm{lo}_k, \mathrm{hi}_k)$ of~$R$.
\end{definition}

\begin{lemma}[The contract of the ranges of a row]\label{lem:siftBandsRanges_spec}
  \lean{JPoints.siftBandsRanges_spec}
  \leanok
  \uses{def:TubeBands.rowRanges, def:bandRowRanges}
  The ranges that \texttt{sift\_bands} computes for the row $(a, b)$ of the band from the row
  $b_{\mathrm{lo}}$ on, with the bands~$\mathrm{bd}$ of the analysis of the tube, are
  $\mathrm{bandRowRanges}(g, \mathrm{sd}, \mathrm{bd}.\mathrm{rowRanges}(b - b_{\mathrm{lo}}), a,
  b, \mathrm{tube})$.
\end{lemma}

\begin{proof}
  \leanok
  \uses{def:rowRegion, def:cutRanges}
  By \texttt{vcgen}, with an invariant of the loop that copies the ranges of the row from the bands
  (\lname{JPoints.CopyRangesInvariant}) and the contract of \texttt{row\_cut} by value, which is
  $\mathrm{cutRanges}$ (\lname{JPoints.rowCut\_spec}).
\end{proof}

\begin{lemma}[The ranges of the sieve by bands]\label{lem:bandRanges_eq_bandRowRanges}
  \lean{JPoints.bandRanges_eq_bandRowRanges}
  \leanok
  \uses{def:bandKeeps, def:TubeBands.rowRanges, def:tubeCell, def:bandOf, def:bandRowRanges, def:Double.ofDyadic}
  If $\mathrm{bandOf}(H, b_{\mathrm{first}}, m, b) = (b_{\mathrm{lo}}, b_{\mathrm{hi}})$, then the
  ranges of the test of the bands (Definition~\ref{def:bandKeeps}),
  $\mathrm{bandRanges}(g, \mathrm{sd}, a, b_{\mathrm{first}}, m, b)$, are the ranges
  $\mathrm{bandRowRanges}$ makes of the ranges of the row
  $b - b_{\mathrm{lo}}$ that the analysis of the band $b_{\mathrm{lo}}, \dots, b_{\mathrm{hi}}$
  collects, in the tube, from the polynomials that \texttt{tube\_a\_init} sets for~$a$ in the
  initial state of the analysis.
\end{lemma}

\begin{proof}
  \leanok
  \uses{lem:siftBandsRanges_spec}
  By the definition of $\mathrm{bandRanges}$ and Lemma~\ref{lem:siftBandsRanges_spec}.
\end{proof}

\begin{lemma}[The points of the ranges]\label{lem:bandRowPoints_eq_rangesPoints}
  \lean{JPoints.bandRowPoints_eq_rangesPoints}
  \leanok
  \uses{def:bandKeeps, def:rowPoints, def:bandRowRanges}
  In the tube or the real region (\texttt{sieve\_mode} $\ne 0$), if the ranges $R =
  \mathrm{bandRanges}(g, \mathrm{sd}, a, b_{\mathrm{first}}, m, b)$ lie in $[w_{\mathrm{low}},
  w_{\mathrm{high}})$ in increasing order, then the points that the sieve by bands finds in the row,
  $\mathrm{bandRowPoints}(g, \mathrm{sd}, a, b_{\mathrm{first}}, m, b)$, are
  $\mathrm{rangesPoints}(g.\mathrm{sift}, \mathrm{sd}, a, b, R)$.
\end{lemma}

\begin{proof}
  \leanok
  \uses{def:orderedTriples, def:HasPointOver, def:colKeeps, def:CheckGlobals, def:rowRangePoints}
  Both list the points over the triples $(a, b, c)$, $c$ increasing, whose third coordinates the
  sieve of the row hands to the exact check off the conic: $\mathrm{bandRowPoints}$ those whose word
  lies in a range of~$R$, $\mathrm{rangesPoints}$ those of each range in turn; as the ranges are in
  increasing order, the two lists agree.
\end{proof}

\begin{lemma}[The contract of \texttt{tube\_word}]\label{lem:tubeWord_spec}
  \lean{JPoints.tubeWord_spec}
  \leanok
  \uses{def:firstIf, def:CheckGlobals, def:Printf.emitPoint, def:IsStageSetup, def:rowRangePoints, def:siftBands}
  Suppose that the statics are set up for the rows of the sieve (Definition~\ref{def:IsStageSetup}),
  that the rows of the sieve tables are those that \texttt{row\_setup} sets for the row $(a, b)$,
  with consistent caches, the index of the current chunk is below the number of the
  chunks, $w_{\mathrm{low}} \le i < w_{\mathrm{high}}$, $0 \le a \le H$, $|b| \le H$, $b \ge 1$ when $a = 0$,
  and the mask and the excluded third coordinate are those of the row
  (\lname{JPoints.TubeWordPre}). Then \texttt{tube\_word} prints the points of the row in its
  word~$i$, $\mathrm{rowRangePoints}(g, \mathrm{sd}, a, b, i, i + 1)$, only the first
  under~\texttt{-1}, says whether it found one under~\texttt{-1}, and changes no field of the state
  but the current chunk, the data of the third stage, the counters and the output.
\end{lemma}

\begin{proof}
  \leanok
  \uses{def:orderedTriples, def:HasPointOver, def:stage1Keeps, def:colKeeps, def:CheckGlobals.WF, def:refIsSquare, def:CurveTables.WF, def:SieveTables.WF, lem:SieveTables.WF.getLsbD_eq_admits12, def:SiftData.WF, lem:checkBits_spec}
  The chunk of the word is found from the current one (\lname{JPoints.seekChunk\_spec}). The mask
  of the row, cut to the row and without the bit of the excluded third coordinate, is and-ed with
  the word of the sieve table of each prime of the first two stages at its residue; an invariant
  over the primes says that the bits left are those of the third coordinates that the tables up to
  the cursor admit (Lemma~\ref{lem:SieveTables.WF.getLsbD_eq_admits12}). The bits left go to
  \texttt{check\_bits} (Lemma~\ref{lem:checkBits_spec}).
\end{proof}

\begin{lemma}[The contracts of the sieve of a range]\label{lem:siftBandsTube_spec}
  \lean{JPoints.siftBandsTube_spec, JPoints.siftBandsRegion_spec}
  \leanok
  \uses{def:firstIf, def:CheckGlobals, def:Printf.emitPoint, def:IsStageSetup, def:rowRangePoints, def:siftBands}
  Suppose the hypotheses of Lemma~\ref{lem:tubeWord_spec} for a range $w_{\mathrm{low}} \le
  \mathrm{lo} < \mathrm{hi} \le w_{\mathrm{high}}$ of the row in place of the word~$i$, with a bit
  array of the size of a chunk (\lname{JPoints.BandRangePre}). Then the sieve of the range in the
  tube, word by word, and in the real region, chunk by chunk, each print the points
  $\mathrm{rowRangePoints}(g, \mathrm{sd}, a, b, \mathrm{lo}, \mathrm{hi})$, only the first
  under~\texttt{-1}, say whether they stopped at one under~\texttt{-1}, and change no field of the
  state but the bit array, the current chunk, the data of the third stage, the counters and the
  output (\lname{JPoints.BandRangePost}).
\end{lemma}

\begin{proof}
  \leanok
  \uses{def:orderedTriples, def:HasPointOver, def:colKeeps, def:CheckGlobals.WF, def:SiftData.WF, def:sift, lem:siftRange_spec, lem:tubeWord_spec}
  An invariant over the words, or over the chunks from that of~$\mathrm{lo}$ on: before an early
  return, the points of the words done are printed (\lname{JPoints.BandRangeCont}). In the tube
  each word is sieved by Lemma~\ref{lem:tubeWord_spec}, in the region the part of the range in each
  chunk by the contract of \texttt{sift\_range} (Lemma~\ref{lem:siftRange_spec}); the points of the
  parts, $c$ increasing, make up those of the range.
\end{proof}

\begin{lemma}[The contract of \texttt{sift\_bands}]\label{lem:siftBands_spec}
  \lean{JPoints.siftBands_spec}
  \leanok
  \uses{def:firstIf, def:rowPoints, def:CheckGlobals, def:Printf.emitPoint, def:SiftState.WF, def:IsSearchSetup, def:siftBands}
  Suppose that the globals and the statics are set up as \texttt{find\_points} sets them up
  ($\mathrm{IsSieveSetup}$, Definition~\ref{def:IsSearchSetup}), the state~$s$ of the thread is
  consistent (Definition~\ref{def:SiftState.WF}), \texttt{sieve\_mode} $\ne 0$, $0 \le a \le H$,
  $m \ge 1$, $b_{\mathrm{first}} \ge -H$, and $b_{\mathrm{first}} \ge 1$ when $a = 0$. Then
  \texttt{sift\_bands} prints the points $\mathrm{bandsPoints}(\mathrm{bg}, \mathrm{sd}, a,
  b_{\mathrm{first}}, m)$ (Definition~\ref{def:rowPoints}), only the first under~\texttt{-1}, says
  whether it stopped at one under~\texttt{-1}, and leaves a consistent state.
\end{lemma}

\begin{proof}
  \leanok
  \uses{def:rowKeeps, def:bandKeeps, def:Cli.Input, def:Cli.Input.WF, def:CheckGlobals.WF, def:IsSearchGlobals, def:Tables.WF, def:SieveTables.WF, def:SiftData.WF, def:IsStageSetup, def:rowRangePoints, lem:rowSetup_spec, def:sift, def:Double.toRat, def:Double.ofDyadic, def:Double.add, def:eval2, def:TubeBands.rowRanges, def:tubeCell, lem:tubeAInit_spec, def:rowRegion, lem:rowRegion_spec, def:cutRanges, def:bandOf, def:bandRowRanges, lem:bandRowPoints_eq_rangesPoints, lem:siftBandsTube_spec, lem:siftBandsRanges_spec}
  Invariants over the bands and over the rows of a band say that, before an early return, the
  points of the rows done are printed. From \texttt{tube\_a\_init} on (Lemma~\ref{lem:tubeAInit_spec})
  the state of the analysis holds the polynomials of~$a$ (\lname{JPoints.TubeState.HoldsPolys}), so
  that the ranges of a row of a band are those of the analysis from the initial state, as
  $\mathrm{bandRanges}$ has them (\lname{JPoints.run\_siftBandsRanges\_siftBandsBand}). A row
  (\lname{JPoints.siftBandsRow\_spec}) that the analysis leaves no word has no ranges, and one that
  the condition at~$2$ drops is not among the rows of $\mathrm{bandsPoints}$; otherwise its ranges,
  in increasing order (Lemma~\ref{lem:rowRegion_spec} and the order of the bands of the tube), are
  sieved one after the other by Lemma~\ref{lem:siftBandsTube_spec} after \texttt{row\_setup}
  (Lemma~\ref{lem:rowSetup_spec}), and their points are those of the row by
  Lemma~\ref{lem:bandRowPoints_eq_rangesPoints}. The rows of the bands in turn are the rows from
  $b_{\mathrm{first}}$ on (\lname{JPoints.rowsBelow\_append\_rowsFrom}).
\end{proof}

\section{Soundness of the test of the bands}\label{sec:bandsSound}

By Lemma~\ref{lem:siftBands_spec}, \texttt{sift\_bands} prints the points over the triples of its
rows that the test of the bands keeps (Definition~\ref{def:bandKeeps}). It remains to show that
this test keeps every triple over which a point lies: in the box it keeps every third coordinate,
in the tube by the analysis of the tube, in the real region by \texttt{row\_region}, and in the
tube with the cut by both.

\begin{lemma}[The real region keeps the points]\label{lem:exists_mem_rangeList_rowRegion_of_points_ne_nil}
  \lean{JPoints.exists_mem_rangeList_rowRegion_of_points_ne_nil}
  \leanok
  \uses{def:orderedTriples, def:HasPointOver, def:CheckGlobals, def:IsSearchSetup, def:NegativeIntervals.Sound, def:rowRegion}
  Suppose that the globals and the statics are set up as \texttt{find\_points} sets them up
  ($\mathrm{IsSieveSetup}$), which gives $H \le 2^{26}$, and that the intervals
  $\mathrm{bg}.\mathrm{neg}$ are sound for the coefficients of the curve searched. Let $(a, b, c)$
  be a triple of the box ($\mathrm{boxTriples}(H)$, Definition~\ref{def:orderedTriples}) off the
  conic $b^2 = 4ac$ with a point over it within the bounds of the search. Then the word
  $\lfloor c/64 \rfloor$ lies in a range of $\mathrm{rowRegion}(\mathrm{bg}.\mathrm{neg},
  \mathrm{sd}, a, b)$.
\end{lemma}

\begin{proof}
  \leanok
  \uses{def:poly, def:kummer, def:Lifts, def:IsNormalized, def:InBounds, def:IsPoint, def:kummerRoots, lem:mem_orderedTriples_iff, def:colKeeps, def:Cli.Input, def:Cli.Input.WF, def:CheckGlobals.WF, def:IsSearchGlobals, def:Tables.WF, def:SieveTables.WF, def:SiftData.WF, def:IsStageSetup, def:RegionExcludes, lem:rowRegion_cover, lem:NegativeIntervals.Sound.not_regionExcludes_of_hasPointOver}
  The word lies in the row. By Lemma~\ref{lem:rowRegion_cover} it lies in a range or is taken out
  for an excluded interval with its margins, and
  Lemma~\ref{lem:NegativeIntervals.Sound.not_regionExcludes_of_hasPointOver} rules out the latter:
  $|a|, |b|, |c| \le H \le 2^{26}$, the polynomial of the set-up is squarefree with at most seven
  coefficients, and $b > 0$ when $a = 0$, since the first nonzero coordinate of a triple of the box
  is positive and the triples $(0, 0, c)$ lie on the conic.
\end{proof}

\begin{lemma}[The tube keeps the points]\label{lem:exists_mem_rowRanges_siftBandsBand_of_points_ne_nil}
  \lean{JPoints.exists_mem_rowRanges_siftBandsBand_of_points_ne_nil}
  \leanok
  \uses{def:orderedTriples, def:HasPointOver, def:CheckGlobals, def:IsSearchSetup, def:TubeBands.rowRanges, def:tubeCell, def:Double.ofDyadic}
  Suppose that the globals and the statics are set up as \texttt{find\_points} sets them up, that
  the fourth coordinate is bounded, that $-H \le b_{\mathrm{lo}} \le b \le b_{\mathrm{hi}} \le H$
  with $b_{\mathrm{hi}} - b_{\mathrm{lo}} + 1 \le \mathtt{TUBE\_BAND}$, and that $(a, b, c)$ is a
  triple of the box with a point over it within the bounds of the search. Then the word
  $\lfloor c/64 \rfloor$ lies in a range of the row $b - b_{\mathrm{lo}}$ of the bands that the
  analysis of the band $b_{\mathrm{lo}}, \dots, b_{\mathrm{hi}}$ in the tube collects from the
  polynomials that \texttt{tube\_a\_init} sets for~$a$ in the initial state of the analysis.
\end{lemma}

\begin{proof}
  \leanok
  \uses{lem:mem_orderedTriples_iff, def:colKeeps, def:Cli.Input, def:Cli.Input.WF, def:CheckGlobals.WF, def:kummerInit, def:IsSearchGlobals, def:Tables.WF, def:SieveTables.WF, def:SiftData.WF, def:IsStageSetup, def:TubeSetup, def:TubeRoot, lem:CheckGlobals.WF.tubeRoot_of_points_ne_nil, def:CellCovers, lem:TubeSetup.cellCovers_run_tubeCell, lem:tubeAInit_spec}
  The point gives the Kummer quadratic of $(a, b, c)$ a root~$d$ with $|d| \le W$
  (Lemma~\ref{lem:CheckGlobals.WF.tubeRoot_of_points_ne_nil}), which the tests of
  \texttt{tube\_cell} on the cells that hold $(b, c)$ do not rule out, so that the analysis collects
  the word (Lemma~\ref{lem:TubeSetup.cellCovers_run_tubeCell}); where \texttt{tube\_ok} fails, the
  band is collected whole.
\end{proof}

\begin{lemma}[The cut keeps the words of both]\label{lem:exists_mem_cutRanges}
  \lean{JPoints.exists_mem_cutRanges}
  \leanok
  \uses{def:cutRanges}
  If the ranges of the lists $\mathit{rs}$ and~$\mathit{gs}$ are each in increasing order ($r_2 \le
  s_1$ for a range~$r$ before a range~$s$) and the word~$w$ lies in a range of~$\mathit{rs}$ and in
  a range of~$\mathit{gs}$, then~$w$ lies in a range of $\mathrm{cutRanges}(\mathit{rs},
  \mathit{gs})$.
\end{lemma}

\begin{proof}
  \leanok
  By the induction of $\mathrm{cutRanges}$: the cut keeps the intersection of the two ranges that it
  compares, and it leaves a range behind only when the next ranges of the other list begin after
  its end.
\end{proof}

\begin{lemma}[Soundness of the test of the bands]\label{lem:bandKeeps_of_points_ne_nil}
  \lean{JPoints.bandKeeps_of_points_ne_nil}
  \leanok
  \uses{def:orderedTriples, def:HasPointOver, def:bandKeeps, def:CheckGlobals, def:IsSearchSetup}
  Suppose that the globals and the statics are set up as \texttt{find\_points} sets them up, that
  $m \ge 1$, $-H \le b_{\mathrm{first}} \le b$ and $m \mid b - b_{\mathrm{first}}$, and that
  $(a, b, c)$ is a triple of the box off the conic with a point over it within the bounds of the
  search. Then the test of the bands keeps it:
  $\mathrm{bandKeeps}(\mathrm{bg}, \mathrm{sd}, a, b_{\mathrm{first}}, m, b, c)$.
\end{lemma}

\begin{proof}
  \leanok
  \uses{lem:mem_orderedTriples_iff, def:Cli.Input, def:Cli.Input.WF, def:CheckGlobals.WF, def:IsSearchGlobals, def:Tables.WF, def:SieveTables.WF, def:SiftData.WF, def:IsStageSetup, def:NegativeIntervals.Sound, def:TubeBands.rowRanges, def:tubeCell, lem:tubeAInit_spec, def:rowRegion, lem:rowRegion_spec, lem:exists_mem_rangeList_rowRegion_of_points_ne_nil, def:cutRanges, lem:exists_mem_cutRanges, lem:exists_mem_rowRanges_siftBandsBand_of_points_ne_nil, def:bandOf, def:bandRowRanges, lem:bandRanges_eq_bandRowRanges}
  In the box the test keeps every third coordinate. Otherwise the ranges of the test are those of
  the band $\mathrm{bandOf}(H, b_{\mathrm{first}}, m, b)$
  (Lemma~\ref{lem:bandRanges_eq_bandRowRanges}), a band of at most \texttt{TUBE\_BAND} rows within
  $[-H, H]$ that holds~$b$. In the tube the fourth coordinate is bounded, and
  Lemma~\ref{lem:exists_mem_rowRanges_siftBandsBand_of_points_ne_nil} applies. In the tube or the
  real region the third stage of the sieve has primes, so that the set-up, with Sturm's theorem as
  the main theorems assume it, makes the intervals on which~$f$ is negative sound, and Lemma~\ref{lem:exists_mem_rangeList_rowRegion_of_points_ne_nil} applies. With the cut,
  both apply, the ranges of the tube and of the real region are each in increasing order
  (Lemma~\ref{lem:rowRegion_spec}), and Lemma~\ref{lem:exists_mem_cutRanges} gives the claim.
\end{proof}

With the soundness of the other tests (Chapter~\ref{chap:theorems}), this shows that the filtered
enumeration loses no point in the rows that \texttt{sift\_bands} sieves: the points that it prints
by Lemma~\ref{lem:siftBands_spec} are all the points of the specification over its rows.
